\chapter{Growth Operators and the Algebra $\Dop$}\label{ch:growth}

\begin{intu}\Intu
Learning \emph{is} change of shape. The complex binds organs that work
together, prunes idle ones, and --- when a procedure repeatedly succeeds
--- promotes it to a first-class operator. Four axes of growth, one of
which grows the very set of moves the system can make.
\end{intu}

\section{The structural operators}

\begin{definition}[Base moves]\label{def:ops}
The base moves on $\CC$ are the typed operators
\[
  \{\opbind,\ \opcollapse,\ \opsplit,\ \opmerge\},
\]
each logged to an immutable mutation history.
\end{definition}

\begin{remark}[``Pachner move'' was discarded, correctly]
Earlier drafts called these Pachner moves. A knowledge complex is not a
manifold, so the invariance guarantees that make Pachner moves useful
--- that any two triangulations of a PL manifold are connected by a
finite sequence of them --- simply do not transfer. Borrowing the name
would have imported a theorem that does not apply.
\end{remark}

\begin{construction}[The four growth axes]\label{con:axes}
Evolution proceeds along
\[
  \underbrace{\delta\CC_{\text{fine}}}_{\text{concepts}},\quad
  \underbrace{\delta\CC_{\text{coarse}}}_{\text{coalitions}},\quad
  \underbrace{\delta\Fs}_{\text{stalk and restriction data}},\quad
  \underbrace{\delta\Dop}_{\text{new operators}} .
\]
The fourth axis is the essential one: the operator algebra is
time-dependent, $\Dop(t)\to\Dop(t+1)$.
\end{construction}

\begin{definition}[Hebbian coupling and decay]\label{def:hebb}
For a coalition $\sigma$ with weight $w(\sigma,t)\in[0,1]$:
co-activation increments $w$; idle steps apply
$w(\sigma,t+1)=(1-\gamma)w(\sigma,t)$ for a decay rate $\gamma\in(0,1)$;
a coalition falling below a bind threshold \emph{collapses}. A pair past
the bind threshold is a $1$-simplex of $\Kc$.
\end{definition}

\begin{proposition}[Idle coalitions collapse in finite time]\label{prop:decay}
\Proved\ If $\sigma$ is idle from time $t_0$ and the bind threshold is
$\theta>0$, then $\sigma$ collapses at the first
$t\ge t_0+\log_{1/(1-\gamma)}\bigl(w(\sigma,t_0)/\theta\bigr)$.
\end{proposition}
\begin{proof}
Iterating the decay gives $w(\sigma,t_0+m)=(1-\gamma)^mw(\sigma,t_0)$,
which falls below $\theta$ as soon as
$m\ln(1-\gamma)<\ln\bigl(\theta/w(\sigma,t_0)\bigr)$; since
$\ln(1-\gamma)<0$ this is $m>\log_{1/(1-\gamma)}(w(\sigma,t_0)/\theta)$.
\end{proof}

\begin{construction}[Consolidation and microneurons]\label{con:micro}
Over the weight filtration, persistent homology reads which coalitions
and concept-cycles \emph{survive}; survivors are promoted into permanent
structure. A \emph{microneuron} is a promoted sub-tree: when a composite
DAG repeatedly resolves conflict \emph{and} passes verification, it is
named and thereafter emitted by the planner as a single atomic generator
of $\Dop$. This is $\delta\Dop$ made concrete. The promotion gate is
\textsc{VerifyOp} (Chapter~\ref{ch:verify}): a composite becomes a
generator only if it survives the external oracle.
\end{construction}

\begin{hon}\Hon
The requirement that promotion pass \textsc{VerifyOp} is doing real
work. Without it, ``repeatedly resolved conflict'' promotes whatever
procedure most reliably makes the organs agree --- which, in a system
whose only cheap signal is agreement, selects for procedures that
manufacture consensus rather than truth. Coherence is a control signal,
not a fitness function, and the gate is what keeps that distinction
enforced.
\end{hon}

\chapter{The Fine Dynamics}\label{ch:flow}

\begin{intu}\Intu
Inside an organ, reasoning is a \emph{flow}. New material perturbs the
internal state and a gradient descent pulls it back toward internal
coherence, subject to hard constraints (axioms that must hold) and soft
constraints (inferred facts, bendable).
\end{intu}

\begin{construction}[Fine state and deductive flow]\label{con:flow}
The fine state of an organ is a configuration $x\in(\Rd)^m$ over its $m$
active concepts, with a scalar \emph{conflict functional} $\Phi(x)\ge0$
measuring internal inconsistency --- the fine-level analogue of
$\discord$. Reasoning integrates
\begin{equation}\label{eq:flow}
  \dot x=-\nabla\Phi(x)
\end{equation}
by a fixed-step fourth-order Runge--Kutta scheme. The reported
\emph{$\delta$-strain} of a step is the achieved change $\Delta\Phi$.
Two constraint types shape $\Phi$: \emph{rigid} constraints
(user-stated axioms) act as hard orthogonal projections onto an affine
subspace via a projector $P$; \emph{fluid} constraints (inferred facts)
act as soft quadratic penalties.
\end{construction}

\begin{proposition}[The flow is a descent]\label{prop:descent}
\Proved\ Along any solution of \eqref{eq:flow},
$\frac{d}{dt}\Phi(x(t))=-\norm{\nabla\Phi(x(t))}^2\le0$, with equality
exactly at critical points. If $\Phi$ is bounded below --- which it is,
by $\Phi\ge0$ --- then $\Phi(x(t))$ converges and
$\nabla\Phi(x(t))\to0$ along a subsequence.
\end{proposition}
\begin{proof}
Chain rule: $\frac{d}{dt}\Phi=\inner{\nabla\Phi}{\dot x}
=-\inner{\nabla\Phi}{\nabla\Phi}$. Monotone and bounded below implies
convergent; integrating
$\int_0^\infty\norm{\nabla\Phi}^2\,dt=\Phi(x(0))-\lim\Phi<\infty$
forces $\liminf\norm{\nabla\Phi}=0$.
\end{proof}

\begin{remark}[Observed sign behaviour]
On a live run, injecting a new concept raised $\Phi$ by a structural
jump ($\delta$-strain $=+0.202$) and the subsequent flow pulled it back
down ($\delta$-strain $=-0.081$): a structural expansion followed by a
deductive relaxation, which is the intended dynamics. Note that the
expansion is \emph{not} a descent step --- it is an exogenous change of
the functional itself, and only the relaxation is governed by
Proposition~\ref{prop:descent}.
\end{remark}

\begin{hon}\Hon
The closed form of $\Phi$ is an implementation-level choice --- a
penalised quadratic on the fine complex --- and is deliberately not
over-specified here. What is mathematically committed is \eqref{eq:flow}
(a gradient flow on a nonnegative conflict functional) and the
rigid/fluid split as hard projection versus soft penalty. Convergence to
a \emph{global} minimum is not claimed and would be false for a
non-convex $\Phi$.
\end{hon}

\chapter{Scheduling: Antichains, and Why ``Operad'' Is Not Earned}\label{ch:operad}

\begin{intu}\Intu
A reasoning act is a small program: a directed acyclic graph of
operators, some depending on others. Independent operators can run at
once, and the scheduler slices the DAG into layers of mutually
independent work.
\end{intu}

\begin{definition}[Antichain foliation]\label{def:foliation}
A reasoning act is a DAG $T=(N,\prec)$ of operator nodes. The scheduler
foliates $T$ into the antichain layers of $\prec$: layer $0$ is the set
of minimal nodes, and layer $\ell+1$ is the set of nodes all of whose
predecessors lie in layers $\le\ell$. Each layer is executed in
parallel. Correctness follows because nodes in one antichain share no
$\prec$-dependency. Node \emph{support} sets index which fine-complex
vertices an operator reads and writes, so slices with disjoint support
are safe to co-execute.
\end{definition}

\begin{remark}[The intended structure is a trace monoid]\label{rem:tracemonoid}
What Definition~\ref{def:foliation} describes is antichain layering of a
DAG. An operad requires multi-input composition together with
equivariance, associativity and unit axioms, none of which are stated or
used anywhere. The structure the design is reaching for is a
\emph{trace monoid} $\mathbb M(\Sigma,I)$ in the sense of Mazurkiewicz:
the free monoid on the operator alphabet modulo $ab=ba$ for pairs in an
\emph{independence relation} $I$, whose canonical Foata normal form is
exactly the antichain layering the scheduler computes. That
identification would be a \emph{sharpening} rather than a downgrade,
since it supplies normal forms, a Hilbert series, and a canonical form
for deduplication. The terminology should be corrected regardless of
whether the identification is completed.
\end{remark}

\begin{hon}\Hon
\Refuted\ The trace-monoid identification \textbf{does not currently
hold}, and the reason is three separate defects, measured 2026-07-29
against the deployed operator set.

\emph{(i) The relation is not symmetric.} In the scheduler, admitting
\emph{any} empty-support node --- including a read-only one --- marks
the whole slice as a global mutation and blocks every later node. So $8$
of the $15$ operator pairs can share a slice in one order and not in the
other. An independence relation is symmetric by definition; a
non-symmetric relation defines no trace monoid, hence no Foata normal
form and no Hilbert series. This is very likely an unintended scheduling
behaviour and is worth fixing on its own merits.

\emph{(ii) The support sets are placeholders.} Every support accessor
returns a hardcoded literal --- $\varnothing$, $\{0\}$ or $\{1\}$ ---
with three of the six operators returning the same $\{0\}$, and a
comment standing in for an analysis of what the operator actually
touches. Any invariant computed from these describes six constants, not
the system.

\emph{(iii) The convention for $\varnothing$ is inverted.} In the engine
an empty support means \textsc{global}: touches everything, must run
isolated. Read as a set, $\varnothing$ is disjoint from everything and
would mean independent of everything. The two readings give capacities
differing by a factor of $1.8$ ($\kappa=3.00$ versus $\kappa=5.45$,
against $\kappa=6$ for the free monoid), so the quantity is not well
defined until the convention is fixed.

Under the reading the engine actually implements, the algebra is
\emph{nearly free} --- about $9\%$ below the free monoid --- because
three of six operators share one support. The order of repair is: fix
the symmetry defect; derive real read/write sets, including the shared
state the current model ignores entirely (the store, the obstruction
counter $\Omega$, the mutation log, and the coarse stalks that $\pi_v$
overwrites after every execution); only then recompute $I$. Operators
that all mutate $\Omega$ do not commute whatever their vertex supports
say. Until that is done, no capacity, Koszulity, or resolution claim
about $\Dop$ should appear anywhere.
\end{hon}

\chapter{The Lift--Project Adjunction}\label{ch:adjunction}

\begin{intu}\Intu
Two worlds: free-form natural language, and structured reasoning DAGs.
One map \emph{lifts} language into a validated DAG; a partner map
\emph{projects} a DAG back into language. Meaning that carries no
reasoning content collapses under the lift --- it lands in the kernel.
\end{intu}

\begin{construction}[The boundary morphisms]\label{con:adjunction}
Let $\mathcal L$ send a natural-language request to a validated
reasoning DAG, and $\mathcal M$ send a DAG to a natural-language
rendering. The lift is guarded by a validator rejecting malformed DAGs
at the boundary: duplicate identifiers, dangling children, cycles by
depth-first colouring, absence of a terminal output node, and
disconnection. The downstream engine therefore only ever receives
well-formed programs.
\end{construction}

\begin{remark}[The kernel behaves as intended]
Semantically void input maps into $\ker\mathcal L$: a joke embedded in a
mathematical prompt produced zero DAG nodes, exactly as a kernel element
should. This is an observation on one input, not a theorem.
\end{remark}

\begin{hon}\Hon
$\mathcal L\dashv\mathcal M$ is the intended categorical \emph{framing}
--- lift left adjoint to project --- and it organises the design
usefully. The unit and counit are not constructed here and the triangle
identities are not verified, so ``adjunction'' should be read as a
structuring principle plus two concrete morphisms with an observed
kernel behaviour, and not as a proved theorem. Stating this is part of
the honesty invariant rather than an apology for it.
\end{hon}
